%@ pages 447-448
%@ status translated
% Bibliography (pp. 447-448).
%
% Convention: titles are given in English, with the title as printed in the
% original in parentheses.  This file is excluded from the French-leftover
% check in tools/translation_report.py, since French, German and Japanese
% titles are expected here.

\begin{thebibliography}{20}

\bibitem{auslander-goldman}
Auslander (M.) and Goldman (O.). \emph{The Brauer group of a commutative
ring}. Trans. Amer. Math. Soc., t.~97, 1960, p.~367--409.

\bibitem{azumaya-1}
Azumaya (Gorô). \emph{Corrections and supplementaries to my paper concerning
Krull--Remak--Schmidt's theorem}. Nagoya Math. J., t.~1, 1950, p.~117--124.

\bibitem{azumaya-2}
Azumaya (Gorô). \emph{On maximally central algebras}. Nagoya Math. J., t.~2,
1951, p.~119--150.

\bibitem{bourbaki-algebre}
Bourbaki (Nicolas). \emph{Algebre}, Chapitre~8: \emph{Modules et anneaux
semi-simples} (Modules and semisimple rings). Paris, Hermann, 1958
(Act. Scient. et Ind., 1261; \emph{Éléments de Mathématique}, 1).

\bibitem{bourbaki-commutative}
Bourbaki (Nicolas). \emph{Algebre commutative} (Commutative algebra) (to
appear).

\bibitem{cartan-eilenberg}
Cartan (H.) and Eilenberg (S.). \emph{Homological algebra}. Princeton,
Princeton University Press, 1956 (Princeton Mathematical Series, 19).

\bibitem{chevalley-grothendieck}
Chevalley (C.) and Grothendieck (A.). \emph{Catégories et foncteurs}
(Categories and functors) (to appear).

\bibitem{eckmann-schopf}
Eckmann (B.) and Schopf (A.). \emph{Über injektive Moduln} (On injective
modules). Archiv der Math., t.~4, 1953, p.~75--78.

\bibitem{goldie}
Goldie (A.~W.). \emph{Semi-prime rings with maximum condition}. Proc. London
Math. Soc., Series~3, t.~10, 1960, p.~201--220.

\bibitem{grothendieck-1}
Grothendieck (Alexander). \emph{Sur quelques points d'algèbre homologique}
(On some points of homological algebra). Tôhoku Math. J., Série~2, t.~9,
1957, p.~119--221.

\bibitem{grothendieck-2}
Grothendieck (Alexander). \emph{Technique de descente et théorèmes
d'existence en géométrie algébrique}, II: \emph{Le théorème d'existence en
théorie formelle des modules} (Descent technique and existence theorems in
algebraic geometry, II: The existence theorem in the formal theory of
modules). Séminaire Bourbaki, t.~12, 1959--1960,
n\textsuperscript{o}~195, 22 pages.

\bibitem{grothendieck-dieudonne}
Grothendieck (A.) and Dieudonné (J.). \emph{Éléments de géométrie
algébrique}, I: \emph{Le langage des schémas} (Elements of algebraic
geometry, I: The language of schemes). Paris, Presses universitaires de
France, 1960 (Institut des Hautes Études Scientifiques, Publications
mathématiques, 4).

\bibitem{leptin-1}
Leptin (Horst). \emph{Linear kompakte Moduln und Ringe}, I. Math. Z.,
t.~62, 1955, p.~241--267.

\bibitem{leptin-2}
Leptin (Horst). \emph{Linear kompakte Moduln und Ringe}, II. Math. Z.,
t.~66, 1956--1957, p.~289--327.

\bibitem{lesieur-croisot-1}
Lesieur (L.) and Croisot (R.). \emph{Structure des anneaux premiers
nothériens} (Structure of noetherian prime rings). C.~R. Acad. Sc. Paris,
t.~248, 1959, p.~2545--2547.

\bibitem{lesieur-croisot-2}
Lesieur (L.) and Croisot (R.). \emph{Théorie nothérienne des anneaux, des
demi-groupes et des modules dans le cas non commutatif}, II. Math. Annalen,
t.~134, 1957--1958, p.~458--476.

\bibitem{matlis}
Matlis (Eben). \emph{Injective modules over noetherian rings}. Pacific
J.~of Math., t.~8, 1958, p.~511--528.

\bibitem{morita}
Morita (Kiiti). \emph{Duality for modules and its applications to the theory
of rings with minimum condition}. Sc. Rep. Tokyo Kyoiku Daigaku, t.~6,
1958--1959, p.~83--142.

\bibitem{serre-groupes}
Serre (Jean-Pierre). \emph{Groupes proalgébriques} (Proalgebraic groups).
Paris, Presses universitaires de France, 1960 (Institut des Hautes Études
Scientifiques, Publications mathématiques, 7).

\bibitem{shih}
Shih (Weishu). \emph{Ensembles simpliciaux et opérations cohomologiques}
(Simplicial sets and cohomological operations). Séminaire Cartan,
t.~11, 1958--1959: \emph{Invariant de Hopf et opérations cohomologiques}
(Hopf invariant and cohomological operations), n\textsuperscript{o}~7,
10 pages.

\end{thebibliography}